\documentclass[11pt]{article}
\usepackage[a4paper,margin=25mm]{geometry}
\usepackage{amsmath,amssymb,amsthm,hyperref}
\newtheorem{theorem}{Theorem}\newtheorem{lemma}{Lemma}
\title{Three simultaneous polynomial-size dice in fully fair families}
\author{Complete proof; three independent reviews passed}\date{30 September 2026}
\begin{document}\maketitle
\begin{theorem}\label{main}
There are absolute constants $C,c>0$ such that for every sufficiently
large $n$ there is a fully permutation-fair family of $n$ finite dice
in which three specified dice have at most $C n^c$ faces each.
All faces on every die are equiprobable and all labels are globally
distinct. The other $n-3$ dice may have one common finite face count;
no polynomial bound on that common count is asserted.
\end{theorem}
The exponent is deliberately not optimized. The proof controls one
common denominator for all initial cell means, then performs one final
local insertion. It does not by itself establish a statement for four
or for an arbitrary fixed number of small dice.

\begin{lemma}[One further variable from controlled cell means]\label{insert}
Suppose rational densities $f_1,f_2$ are supported on complementary
pure cells forming a rational partition of $[0,1]$. Assume:
\begin{enumerate}
\item the number of cells, inverse cell lengths, inverse positive cell
masses, and polynomial piece degrees are bounded by a polynomial in $m$;
\item the conditional density on each cell, rescaled to $[0,1]$, is
bounded above and below by absolute positive constants;
\item there is one polynomial integer $D$ such that every endpoint,
every positive cell mass, and every unnormalized first cell moment
belongs to $D^{-1}\mathbb Z$.
\end{enumerate}
Then there are three rational piecewise polynomial densities
$\widetilde f_1,\widetilde f_2,\widetilde f_3$ on pairwise disjoint
pure cells whose ordered polynomial moments through total degree $m$
coincide with those of independent variables having densities
$f_1,f_2,1$. Their cell masses belong respectively to
$N_1^{-1}\mathbb Z,N_2^{-1}\mathbb Z,N_3^{-1}\mathbb Z$ for
polynomial integers $N_i$. Their rescaled conditional density bounds
are absolute, and their polynomial piece degrees are polynomial in $m$.
There is no required bound on the new cell endpoint or higher-moment
denominators, nor on the number of new cells beyond finiteness.
\end{lemma}
\begin{proof}
Consider an original cell $I=[a,b]$, length $\ell$, occupied by old
label $i\in\{1,2\}$. Put
\[
 \gamma=\int_I f_i(t)dt,\qquad z=\int_I t f_i(t)dt.
\]
After rescaling $I$ to $[0,1]$, its old conditional density is a rational
polynomial $w$ with mean
\[
 \mu=\frac{z-a\gamma}{\gamma\ell}.
\]
Choose polynomial integers $N_1=N_2=N$ and $N_3=T$, both multiples
of $D$, sufficiently large, and prescribe the local counts
\[
 K_I=N\gamma,\qquad L_I=T\ell.
\]
These are positive integers. Moreover
\begin{equation}\label{integrality}
 K_I L_I\mu=NT(z-a\gamma)\in\mathbb Z,
\end{equation}
since $D^2(z-a\gamma)$ is integral.

Apply the theorem in
\path{../size_bounds/generalized_local_two_variable_completion.tex}
on every cell, with its old density $w$ as the theorem's $B$ variable
and its uniform variable as the theorem's $A$ variable. It is applicable
for all integers
\[
 K_I\ge C(m+\deg w+2)^4,\qquad
 L_I\ge C[K_I(m+\deg w+2)]^{160}
\]
satisfying~\eqref{integrality}. First take $N$, a multiple of $D$,
as a sufficiently large fixed polynomial in $m$ to satisfy every
first inequality. Then take $T$, another multiple of $D$, as a
sufficiently large fixed polynomial to satisfy every second inequality.
The uniform bounds in the hypotheses, the polynomial lower bounds on
$\gamma,\ell$, and the common $D$ make these simultaneous choices
possible. No least common multiple of independently produced cell
denominators is taken.

Rescale each local output back to $I$. Multiply the old density by its
original cell mass $\gamma$, and the new density by the original
uniform mass $\ell$. Each local old box has global mass
$\gamma/K_I=1/N$. Each local new gap has mass an integer times
$\ell/L_I=1/T$. Assemble the old pieces over their respective original
cells to obtain $\widetilde f_1,\widetilde f_2$, and assemble the new
pieces over all original cells to obtain $\widetilde f_3$.
Each integrates to one. Their supports are pairwise pure, every
coefficient and endpoint is rational, and their claimed common mass
denominators are $N,N,T$. The local theorem gives the uniform
conditional density bounds and polynomial piece degrees.

It remains to check the global moments. Partition an ordered integral
according to the three original cells occupied by its variables.
The first two variables never occupy the same original cell.
If all three cells are distinct, their relative order is fixed and
the integrand is an ordinary polynomial; the local marginal moment
identities preserve it. If the third variable shares one of the first
two cells, the ordering constraint on that pair is the only nontrivial
local chamber condition. The local two-variable theorem preserves
all its polynomial moments through degree $m$, and the other old
variable is in a different cell with fixed relative order. Expand
the polynomial in the other variable and use its preserved marginal
moments. These cases exhaust the partition. The original cell masses
are unchanged, so summing proves every asserted ordered moment identity.
\end{proof}

\begin{proof}[Proof of Theorem~\ref{main}]
Write $m=n-3$. Use
\path{two_pure_densities_with_rational_means.tex}
with parameter $m+1$ to obtain $f_1,f_2$ such that these variables
and $m+1$ independent uniforms are fully permutation-fair. Its cell
bounds and single polynomial denominator verify every hypothesis of
Lemma~\ref{insert}. Separate one of those uniforms as the third
specified variable and apply that lemma through degree $m$.
Conditional on the values of the three specified variables, a fixed
full order of the remaining $m$ uniforms has probability, on its
specified chamber,
\[
 \prod_{j=0}^3\frac{(t_{j+1}-t_j)^{r_j}}{r_j!},
 \qquad t_0=0,\ t_4=1,\ \sum_jr_j=m.
\]
This is a polynomial of total degree $m$. The lemma preserves its
integral in each chamber. Thus the three new rational pure densities
and $m$ uniforms are still fully permutation-fair.

Let their polynomial mass denominators be $N_1,N_2,N_3$. Choose a
common polynomial integer $H$ sufficiently large. On each pure cell
belonging to label $i$, its mass is an integer $a/N_i$. Replace its
conditional density by an equal-weight real quadrature using exactly
$Ha$ nodes, exact through degree at least $2m+2$.
Every sufficiently large number of nodes is permitted, uniformly
in the interval and cell: the separated quadrature lemma in
\path{../size_bounds/generalized_local_two_variable_completion.tex}
applies after rescaling. Its hypotheses are exactly the absolute upper
and lower density bounds and the polynomial degree bound established
above. Taking its exactness degree to be the larger of $2m+2$ and one
more than that piece degree, one polynomial choice of $H$ works on
all cells. It gives interior, pairwise distinct nodes directly.
Increase $H$ also to ensure $H\ge m^2$.

The resulting specified dice have exactly $HN_i$ equiprobable nodes,
polynomial in $m$. On each product of pure cells their relative order
is fixed, so replacing each density by its conditional quadrature
preserves every degree-$m$ order kernel. They remain fully fair with
the $m$ uniforms. Each pure cell supplies a cluster containing at
least $m^2$ distinct nodes. Shrink each cell to a closed subinterval
containing its finitely many nodes in its interior, giving disjoint
cluster closures. The cluster moments through degree $m$ are rational,
because they equal moments of a rational polynomial density on an
interval with rational endpoints, multiplied by its rational cell mass.

All hypotheses of the reviewed finiteization theorem
\path{../size_bounds/clustered_multi_exceptional_lifting.tex}
now hold for $k=3$. It gives a completely finite fair family without
changing the three specified counts $HN_i$. The remaining $m$ dice
have a common finite count. Replacing the global face ordering by
successive integer labels makes all labels distinct. Finally, the three
counts $HN,HN,HT$ can be equalized to $H\operatorname{lcm}(N,T)\le HNT$
by uniform face replication within each specified die. Place each replica
in a sufficiently small disjoint neighborhood of its old label; every
relative order is unchanged. The common specified count is still polynomial.
This proves the theorem.
\end{proof}
\end{document}
